\documentclass{article}
\usepackage[T1]{fontenc}
\usepackage{lmodern}
\usepackage{graphicx}
\usepackage{amsmath,amssymb}
\usepackage[a4paper,margin=2cm]{geometry}
\usepackage[absolute,overlay]{textpos}
\setlength{\TPHorizModule}{1mm}
\setlength{\TPVertModule}{1mm}
\usepackage[indonesian]{babel}
\usepackage{hyperref}
\setlength{\emergencystretch}{2em}
% unit: Halaman 14

\begin{document}

% === Halaman 14 (python/native) ===

14

Contoh 1:  Bob membangkitkan kunci publik dan kunci privatnya. Alice  akan 

mengengkripsi  pesan dengan menggunakan kunci publik Bob.

(a) Pembangkitan kunci (Oleh Bob)


Misal p = 2357, g = 2,  dan x = 1751.


Hitung: y = gx mod p = 21751 mod 2357 = 1185


Hasil:   Kunci publik: (y = 1185, g = 2, p = 2357) 



       Kunci privat: (x = 1751, p = 2357).

      Bob memberitahu kunci publik ini kepada Alice

(b) Enkripsi (Oleh Alice)


Misalkan pesan m = 2035 (nilai m masih berada di dalam selang [0, 2357 – 1]).


Alice memilih bilangan acak k = 1520  (nilai k berada di dalam selang [0, 2357 – 1]). 

syarat g < p , g akar primitif dari 2357 dan 

1  x  p – 2

\end{document}
